\documentclass[11pt]{article}

\usepackage{amsmath,amssymb}
\usepackage{xurl}
\usepackage[hidelinks]{hyperref}
\setlength{\emergencystretch}{2em}

% Shared commands for notes in docs/paper-gaps/.
% Notation follows the LDT paper (references/ldt-paper/).

\newcommand{\Lean}{\textsc{Lean}}
\newcommand{\leanid}[1]{\textsf{#1}}
\newcommand{\ghissue}[1]{\href{https://github.com/LionSR/MIPStarRE/issues/#1}{GitHub issue \##1}}
\newcommand{\ghpr}[1]{\href{https://github.com/LionSR/MIPStarRE/pull/#1}{GitHub PR \##1}}

% --- Paper-compatible notation ---
\newcommand{\F}{\mathbb{F}}
\newcommand{\N}{\mathbb{N}}
\newcommand{\C}{\mathbb{C}}
\newcommand{\tr}{\operatorname{tr}}
\newcommand{\ot}{\otimes}
\newcommand{\E}{\mathop{\bf E\/}}
\newcommand{\bx}{\mathbf{x}}
\newcommand{\by}{\mathbf{y}}
\newcommand{\bu}{\mathbf{u}}
\newcommand{\bv}{\mathbf{v}}

% Approximation relation (paper uses \approx_\eps and \simeq_\zeta)
\newcommand{\approxeps}{\approx_{\varepsilon}}
\newcommand{\simgeq}{\simeq_{\zeta}}

% Polynomial spaces (paper uses \polyfunc, \polysub, \polymeas)
\newcommand{\polyfunc}[3]{\operatorname{Poly}(#1,#2,#3)}
\newcommand{\polysub}[3]{\operatorname{PolySub}(#1,#2,#3)}
\newcommand{\polymeas}[3]{\operatorname{PolyMeas}(#1,#2,#3)}


\title{The Nonnegativity Hypothesis in the Small-Overlaps Bound}
\author{}
\date{2026-04-30}

\newcommand{\Small}{\mathsf{Small}}

\begin{document}
\maketitle

\section{The assertion}

This note concerns the bound on the sum over small overlaps in the rank-reduction
step of the orthogonalization lemma,
\cite[Section~\emph{Orthogonalization lemma}]{Ji2020LowIndividualDegree}.
\footnote{The discrepancy recorded here was identified in
\href{https://github.com/LionSR/MIPStarRE/issues/938}{GitHub issue \#938}.
In the TeX source consulted here, the estimate occurs in
\path{references/ldt-paper/orthonormalization.tex}, lines 606--612.}
We use the notation of the cited source:
the overlaps $o_{a,i}$ are defined by
\[
  o_{a,i} = \bra{\psi} (\ket{v_{a,i}}\bra{v_{a,i}} \ot I) \ket{\psi},
\]
where $\ket{\psi}$ is the state and each $\ket{v_{a,i}}$ is a unit vector from
the range of a projector. The total number of overlap pairs is $r$, the set
$\Small$ is the complement of the $d$ largest overlaps, and $\zeta$ is the
self-consistency error parameter with $0\le \zeta\le 1/4$.

At this point the paper asserts the estimate
\begin{align}
  \sum_{(a,i)\in\Small} o_{a,i}
    &\le 4\sqrt{\zeta},
      \tag{\path{eq:small-overlaps}}
\end{align}
using the auxiliary bounds
\[
  |\Small| \le 2\sqrt\zeta \cdot r,
  \qquad
  \sum_{a,i} o_{a,i} \le 1+2\sqrt\zeta,
\]
together with the fact that every overlap in $\Small$ is at most every overlap
outside $\Small$. The proof of the auxiliary bounds uses the overlaps
themselves; they are manifestly non-negative because each is an expectation of a
positive operator against a state.

\section{The abstraction in Lean}

The formal development extracts the combinatorial content of the estimate and
proves it for an arbitrary real-valued function, without assuming nonnegativity.
\footnote{The formal lemma is
\leanid{sum\_small\_le\_four\_sqrt} in
\path{MIPStarRE/LDT/MakingMeasurementsProjective/QXPLayer/TruncationCombinatorics.lean},
lines 163--174. It is accompanied by two auxiliary lemmas
\leanid{card\_mul\_sum\_small\_le} (lines 100--119) and
\leanid{card\_univ\_mul\_sum\_compl\_le} (lines 125--148).}

In the statement of the lemma, the overlaps are replaced by a function
$f\colon \alpha\to \R$ on a finite set $\alpha$. The subset $L\subseteq\alpha$
plays the role of the large overlaps (the complement of $\Small$), and the
hypotheses are
\begin{itemize}
  \item $|L|=d$,  $d < |\alpha|$,  and
        $|\alpha| \le (1+2\sqrt\zeta)\,d$;
  \item every value on the complement $L^{\mathsf c}$ is at most every value on
        $L$ (the ordering hypothesis);
  \item $\sum_{x\in\alpha} f(x) \le 1+2\sqrt\zeta$ (the total estimate);
  \item $0\le \zeta\le 1/4$.
\end{itemize}
The conclusion is
\[
  \sum_{s\in L^{\mathsf c}} f(s) \le 4\sqrt\zeta.
\]
The proof does not need the hypothesis that every $f(x)$ is non-negative. It
uses only the ordering hypothesis, the total estimate, and the cardinality
relations; the sign of $f$ plays no role in the calculation.

\section{Comparison}

In the paper, the function $o_{a,i}$ is non-negative by construction, so the
implicit hypothesis $o_{a,i}\ge 0$ is automatically satisfied and the estimate
is valid. The formal lemma is therefore strictly stronger than what the paper
states: it has \emph{fewer} hypotheses (it drops the nonnegativity condition) and
derives the same conclusion.

The difference does not affect soundness. The formal proof of the
orthogonalization lemma instantiates the combinatorial lemma with the actual
overlaps $o_{a,i}$, which are indeed non-negative. In that instantiation the
ordering and total-estimate hypotheses are verified from the properties of the
overlaps, and the lemma then supplies the needed bound on the sum over the small
entries.


\section{Relation with the blueprint}

The blueprint states this lemma as
\path{lem:small-overlap-averaging} in
\cite[Chapter~\emph{Making measurements projective}]{MIPStarREBlueprint}.
\footnote{In the blueprint source consulted here, this is in
\path{blueprint/src/chapter/ch04_projective.tex}, lines 93--106.}
It takes $f:\alpha\to \R$ to be an arbitrary real-valued function and does not
require $f\ge 0$; the hypotheses are exactly the ordering hypothesis, the total
estimate, the size constraints, and $0\le\zeta\le 1/4$.  The blueprint therefore
matches the Lean statement rather than the paper, and the discrepancy with the
paper is already resolved at the blueprint level.

\section{Conclusion}

There is no mathematical discrepancy. The formal lemma is a conservative
strengthening of the argument in the paper. The paper's proof of the
small-overlaps bound uses only the ordering hypothesis, the total estimate, and
the cardinality relations; it never invokes the nonnegativity of the overlaps
except through their definition. The formal development makes this explicit by
removing the nonnegativity condition from the statement. In the application
within the orthogonalization lemma, the stronger statement is specialized back to
the non-negative overlaps, so the formal soundness argument remains faithful to
the paper.

\bibliographystyle{alpha}
\bibliography{docs/paper-gaps/references}

\end{document}